\documentclass[border=0pt]{standalone}
% The shared preamble of every docs diagram (docs/diagrams/<name>.tex).
% dvisvgm's driver first, so patterns and opacity survive into the SVG.
\def\pgfsysdriver{pgfsys-dvisvgm.def}
\usepackage[siunitx,RPvoltages]{circuitikz}
\usepackage{xcolor}
% Latin Modern: outline fonts for every glyph a label uses, µ included.
\usepackage[T1]{fontenc}
\usepackage{lmodern}
\usetikzlibrary{arrows.meta,backgrounds,calc,patterns,positioning}
\tikzset{
  % A white card, so diagrams read in mdBook's dark themes too.
  every picture/.append style={show background rectangle,
    background rectangle/.style={fill=white,inner sep=4pt}},
  % Layout geometry: layers told apart by fill and hatching, never colour alone.
  layer a/.style={draw=black,line width=0.6pt,fill=gray!20},
  layer b/.style={draw=black,line width=0.6pt,fill=gray!45},
  layer c/.style={draw=black,line width=0.6pt,fill=white,postaction={pattern=north east lines}},
  layer d/.style={draw=black,line width=0.6pt,fill=white,postaction={pattern=crosshatch dots}},
  marker/.style={draw=black,dashed,line width=0.6pt},
  % A measured distance and its label.
  dim/.style={{Latex[length=2mm]}-{Latex[length=2mm]},line width=0.5pt},
  dim label/.style={font=\small,fill=white,inner sep=1pt},
  % Where a rule fails: thick, dashed and labelled, not only red.
  violation/.style={draw=red!75!black,line width=1.4pt,dash pattern=on 3pt off 2pt},
  violation label/.style={font=\small\bfseries,text=red!75!black},
  caption/.style={font=\small},
}

\begin{document}
\begin{circuitikz}
  % The inverter.
  \node[pmos] (p) at (0,1.2) {};
  \node[nmos] (n) at (0,-1.2) {};
  \draw (p.D) -- (n.D);
  \draw (p.S) -- ++(0,0.3) node[above] {VDD};
  \draw (n.S) -- ++(0,-0.3) node[below] {VSS};
  \draw (p.G) -- (n.G);
  \draw (0,0) node[circ] {} -- ++(0.6,0) node[right] {Y};
  % A, the inverter's input, pulled to VSS through two series nfets.
  \draw (p.G |- 0,0) node[circ] {} -- (-2.2,0) node[circ] {};
  \node[below] at (-1.55,-0.05) {A};
  \node[nmos] (n1) at (-2.2,-1.0) {};
  \node[nmos] (n2) at (-2.2,-2.8) {};
  \draw (-2.2,0) -- (n1.D);
  \draw (n1.S) -- (n2.D);
  \draw (n2.S) -- ++(0,-0.3) node[below] {VSS};
  \draw (n1.G) -- ++(-0.2,0) node[left] {B};
  \draw (n2.G) -- ++(-0.2,0) node[left] {C};
  % Findings.
  \draw[violation] (-2.1,0.08) -- (-1.05,0.08);
  \node[violation label,anchor=south east,align=right] at (-2.3,0.15) {A: 2 devices to VSS,\\ limit 1};
  \node[violation label,anchor=east,align=right] at (-3.9,-1.9) {B, C: no path\\ at all};
\end{circuitikz}
\end{document}
